\section{Case Study}
\label{sec:case_study}

We demonstrate the effectiveness of \sysname{} through three case studies across different domains: database query optimization, code evolution, and agent memory management. These cases were selected because they exhibit distinct tree evolution characteristics and analytical goals, allowing us to examine how \sysname{} supports temporal overview, multi-granularity structural analysis, and detailed change inspection in practice.

\subsection{Case Study 1: Detecting Loop Optimization in Database Query Optimizer}

A database query optimizer transforms an input SQL query into an efficient execution plan through a sequence of rewrite and optimization operations. In practice, however, interactions among optimization rules may lead to repeated rewrites or even optimization loops, where similar transformations are applied multiple times. Such behaviors are difficult to detect using existing tools, especially when analysts need to manually inspect long optimization traces and compare repeated plan transformations.

\sysname{} helps analysts identify these repeated optimization patterns without requiring strong domain-specific prior knowledge. In the statistics view, the Evolution Overview compresses recurring operation subsequences and exposes loop-like patterns as repeated summarized units, allowing analysts to quickly detect suspicious repetitive behaviors in long optimizer traces. The difference trend line and operation distribution further reveal intervals with repeated bursts of structural activity, helping users localize the suspicious region. After selecting the relevant time interval, the hierarchy view summarizes the repeated plan transformations at multiple levels of granularity, while the tree views expose the detailed node-level edits associated with each iteration. Through this workflow, analysts can recognize that the optimizer repeatedly applies similar rewrites to structurally related subtrees, thereby revealing inefficient or redundant optimization behavior.

\subsection{Case Study 2: Revealing Code Refactoring Patterns through AST Evolution}

Abstract syntax trees (ASTs) provide a hierarchical representation of program structure. During software development, code revisions naturally induce AST evolution, reflecting operations such as function insertion, variable renaming, statement reorganization, and subtree deletion. In our case, each AST contains roughly 100--200 nodes corresponding to Python syntax elements such as \texttt{Module}, \texttt{FunctionDef}, \texttt{Assign}, \texttt{If}, \texttt{For}, \texttt{Return}, and \texttt{Call}. The resulting evolution sequence includes structural modifications such as inserting new functions, renaming variables or functions, moving code blocks, and removing outdated statements.

This case highlights the value of \sysname{} for understanding refactoring patterns in large and partially stable trees. Since most of the AST remains unchanged across adjacent revisions, direct comparison can easily be overwhelmed by unchanged context. \sysname{} addresses this issue by using difference-aware highlighting and multi-granularity summarization. In the tree views, unchanged nodes are visually simplified while changed nodes are emphasized through visual encoding, making local edits easier to recognize. The glyph-based representation of minimum difference forests reduces cognitive load while preserving sufficient structural context for interpretation. Moreover, when analysts hover over a node, unchanged corresponding nodes in neighboring trees are linked by lines, making it possible to visually trace how a node or subtree is preserved, simplified, or reorganized across revisions. This interaction helps reveal refactoring patterns such as repeated local restructuring, systematic renaming, and incremental function decomposition.

\subsection{Case Study 3: Understanding Memory Organization in Agent Memory Trees}

Large language model (LLM) agents often maintain evolving memory structures to store observations, reflections, and higher-level summaries of interaction histories. As the dialogue continues, the memory tree grows and changes through insertion, aggregation, rewriting, and occasional deletion. A key analytical problem is to understand how low-level observations are gradually organized into higher-level reflections and how the memory structure evolves over long interaction sessions.

In this case, we analyze the memory tree of an LLM agent over more than 15 dialogue rounds. The memory structure follows a hierarchical organization similar to prior generative agent frameworks, with high-level categories such as agent persona, reflections, and observations. The dominant evolution pattern is insertion, corresponding to the agent continuously incorporating new information, together with occasional renaming or restructuring when existing memories are refined or reorganized.

\sysname{} helps reveal both the global trend and the local mechanisms of this memory evolution process. In the statistics view, the overview highlights long insertion-dominated intervals while also exposing a few salient transition points with substantial restructuring. The hierarchy view reveals how many low-level observation nodes are progressively grouped into newly formed reflection nodes, making the memory consolidation process visible at an intermediate structural level. In the tree views, analysts can further inspect how multiple observations are absorbed into a new reflection subtree and how the surrounding memory structure is updated accordingly. Through these coordinated views, \sysname{} makes the agent's memory organization process interpretable, showing not only that the tree grows over time, but also how the agent gradually transforms raw observations into more abstract memory structures.

\subsection{Summary}

Across the three cases, \sysname{} supports analysts in identifying salient intervals, detecting repeated or abnormal structural behaviors, and tracing how local changes accumulate into larger evolutionary patterns. The database query optimization case demonstrates the ability of \sysname{} to reveal repeated transformation loops, the AST case shows its value for exposing refactoring patterns in partially stable trees, and the agent memory case highlights its usefulness for understanding hierarchical memory organization over time. Together, these studies suggest that \sysname{} can effectively support tree evolution analysis in diverse domains with different semantic meanings but similar structural dynamics.
